% docs/manual/pdf-header.tex — typography hardening for the PVT Lab Platform
% software manual PDF build. Included via `pandoc -H` (a raw LaTeX header
% inserted into the preamble) from scripts/build_manual.sh.
%
% Addresses table/paragraph overfull \hbox defects (text overflowing into
% neighboring columns or off the page edge) that show up with narrow
% longtable columns and long inline `\texttt{}` code identifiers:
%
%   - \emergencystretch gives the paragraph builder extra stretchable glue
%     to reach for as a last resort before letting a line go overfull.
%   - hyphenat[htt] enables hyphenation inside \texttt (monospace/code)
%     text, which LaTeX disables by default, so long code tokens can break
%     across lines instead of overflowing their (often narrow) table cell.
%   - \tabcolsep is tightened from LaTeX's 6pt default to 4pt, giving every
%     table column a little more usable width.
%   - \sloppy relaxes interword spacing constraints document-wide, trading
%     a little justification tightness for far fewer overfull lines.
%   - fvextra redefines plain \begin{verbatim} (what pandoc emits for a
%     fenced code block with no language tag, e.g. a literal error-message
%     string) to wrap long lines at the page margin instead of running off
%     it — plain LaTeX verbatim text never wraps and untagged fenced
%     blocks don't go through pandoc's highlighted-code path, so this is
%     the one place hyphenat/sloppy can't reach. The text itself (e.g. an
%     exact error string quoted from the app) is untouched.
\emergencystretch=3em
\usepackage[htt]{hyphenat}
\setlength{\tabcolsep}{4pt}
\sloppy
\usepackage{fvextra}
\DefineVerbatimEnvironment{verbatim}{Verbatim}{breaklines=true,breakanywhere=true}
